\subsection{保护基础搜索的限时算法组合}
本轮采用结构初解与自适应有界邻域的组合。对输入图先构造弱连通分量、共享输入消费者、拓扑层次和各 Pipe 工作量索引；给定迁移种子时，必须在场景 B 下重新评价。初始化保留整图候选、原有结构菜单和四次普通联合修复机会，再启动扩展搜索。这里的“四次”只描述初始化阶段，不能把整个新算法误写成固定十二次评价。

令 $f$ 表示重映射、粒度调整、顺序扰动、联合修复、局部插入、拓扑前沿或分量装箱的一类候选。定义最大分量占比
\begin{equation}\label{eq:q2-dominant-ratio}
\chi(G)=\frac{\max_{C\in\operatorname{WCC}(G)}|C|}{|\mathcal O_c|}.
\end{equation}
当 $\chi(G)>0.5$ 时，提高重映射与粒度类的先验份额；存在多个独立分量时才开放分量装箱。这些规则只依赖图结构，不按算例编号查答案。前期优先保留普通联合搜索：前 $32$ 次提议且位于自适应阶段剩余时间的前 $40\%$ 时采用该基础邻域；随后在不同候选类间按收益与耗时调整抽样。令 $G_f$ 为累计相对周期改善、$t_f$ 为该类累计墙钟耗时、$w_f$ 为结构先验，则抽样权重为
\begin{equation}
\widetilde w_f=w_f\left[1+\min\left(8,\frac{100G_f}{\max(0.1,t_f)}\right)\right],\qquad
P(f)=\frac{\widetilde w_f}{\sum_g\widetilde w_g}.
\end{equation}
未获收益的候选类仍保留正权重。由于 $t_f$ 取实测秒数，相同随机种子在不同主机和并发负载下不保证产生完全相同的搜索轨迹；已保存方案的官方重放则与搜索轨迹无关。

每个多核搜索设 $590$ 秒上限，其中预留 $20\%$ 给最终原版评价和物理内存审计，自适应提议至多 $384$ 次。每个提议可以被合法性检查、重复指纹或剩余时间检查提前拒绝，提议数不等于实际官方评价数。停止条件还参考最近四次评价的最大耗时，剩余时间不足其 $1.3$ 倍时不再启动下一项。该经验预留并不是任意图上的严格完成时间证明；正式结果只接纳完整复核成功的输出。

为减少重复生成成本，在同一不可变图对象内，缓存上下文构造、顺序构造和局部代理函数的确定性结果。缓存键包含函数名及完整实参的摘要，返回值反序列化为独立对象，避免后续修改污染缓存；总载荷上限为 $32$ MiB，过大对象绕过缓存。缓存不改变外层随机数调用，不复用旧方案的全局 DDR 时长。划分、映射或顺序变化后，正式评价重新计算完整核内排程、物理内存和全局事件。

候选 $p$ 仅在官方评价合法且 $(T(p),D(p))$ 字典序优于 incumbent 时替换。对实际评价的合法集合 $\mathcal A$，最终返回
\begin{equation}
p^\star\in\arg\min_{p\in\mathcal A}\operatorname{lex}(T^B(p),D(p)).
\end{equation}
这一有限集合择优性质由逐次更新直接归纳得到，既不证明全局最优，也不保证另一台机器在同一秒数内访问同一个 $\mathcal A$。全量实验固定规则、种子与配置，并单独披露历史初解成本。
